\documentclass{article}
\usepackage{microtype}
\usepackage{graphicx}
\usepackage{booktabs}
\usepackage{hyperref}
\usepackage{icml2026}

\begin{document}

\twocolumn[
\icmltitle{Predicting Random-Walk Cover Times on Small Graphs from Cheap Statistics}

\icmlsetsymbol{equal}{*}
\begin{icmlauthorlist}
\icmlauthor{Anonymous Author}{anon}
\end{icmlauthorlist}
\icmlaffiliation{anon}{Anonymous Institution}
\icmlcorrespondingauthor{Anonymous Author}{anon@example.org}
\icmlkeywords{random walks, cover time, graph statistics, empirical study}
\vskip 0.3in
]
\printAffiliationsAndNotice{}

\begin{abstract}
The cover time of a simple random walk is classically bounded in terms of the number of vertices $n$ and edges $m$, for example by $2m(n-1)$. For small graphs these bounds are loose. We ask whether a few easily computed statistics (n, m, diameter, maximum degree) give a more accurate prediction of the cover time for connected graphs with $n\le 30$. We estimate the cover time of 140 synthetic graphs from 15 families by Monte Carlo with 10{,}000 walks per graph (relative standard error $\le 1.2\%$). The universal bound overestimates by a median factor of about 4 (ratio $C/2m(n-1)$ median $0.25$, range $0.005$--$0.59$). A log-linear model in $n$, $m$, diameter and maximum degree predicts unseen families with median absolute relative error 12\% (bootstrap 95\% CI of mean log error $[0.16,0.22]$), versus 65\% for the rescaled bound. The model is reliable on sparse, path-like and random-regular graphs (about 8--13\% error) but fails on complete graphs (median 52\%), grids (35\%) and binary trees (37\%), where the cover-time structure is not captured by these four statistics. These are empirical findings on a small, synthetic test set, not a general theorem.
\end{abstract}

\section{Introduction}
For a connected graph $G$ with $n$ vertices and $m$ edges, the cover time $C(G)$ is the expected number of steps a simple random walk needs to visit every vertex. Aleliunas et al.~\cite{aleliunas1979random} showed $C(G)\le 2m(n-1)$ for every start vertex, and Feige~\cite{feige1995tight} showed that the worst case over graphs is $\Theta(n^3)$ with a sharp constant, attained by lollipop-type graphs. These are worst-case results. For a specific graph with $n\le 30$, they can be off by an order of magnitude, and a practitioner who wants a quick estimate of $C(G)$ has no standard tool beyond the bound (see the survey by Lov\'asz~\cite{lovasz1993random}).

We ask a narrow empirical question: for connected graphs with at most 30 vertices, do a few statistics that are trivial to compute (vertex count, edge count, diameter, maximum degree) predict $C(G)$ much better than the universal bound, and on which kinds of graphs does such a prediction break down?

Our contributions are: (i) a Monte Carlo dataset of cover-time estimates for 140 graphs in 15 families, with uncertainty; (ii) a comparison of a rescaled universal bound against log-linear predictors, evaluated by leave-one-family-out cross-validation, which tests generalisation to unseen graph types; and (iii) a per-family analysis of where the predictor fails.

\section{Method}
\paragraph{Graphs.} We generate connected graphs with $6\le n\le 30$ at $n\in\{6,10,14,18,22,26,30\}$ in 15 families: paths, cycles, stars, complete graphs, wheels, heap-shaped binary trees, random labelled trees (3 seeds per $n$), random 3- and 4-regular graphs, lollipops, barbells, $r\times c$ grids with $rc=n$ (when $r>1$), and Erd\H{o}s--R\'enyi graphs $G(n,p)$ with $p\in\{0.1,0.2,0.4\}$ conditioned on connectivity (2 seeds each). Each family contributes 7 graphs except random trees (21), and the $G(n,p)$ families (14 each). The number of graphs in each family is small, so per-family numbers should be read as indicative.

\paragraph{Cover time.} For each graph we simulate $10{,}000$ independent walks, each started at a uniformly random vertex, and record the first time all vertices are visited. The sample mean $\hat C$ and its standard error are reported. Starting from a uniform vertex (rather than the worst vertex) makes the target a single well-defined number per graph; it is the quantity that a user would typically want when a start vertex is not specified. The simulator was checked against exact values: for $C_{10}$ the estimate is $45.4\pm0.4$ versus the exact $45$, and for $K_{10}$ it is $25.2\pm0.15$ versus the exact $(n-1)H_{n-1}=25.46$ (both within about two standard errors). Across all graphs the median relative standard error is 0.54\% (maximum 1.15\%), far smaller than the model errors reported below.

\paragraph{Predictors.} All models predict $\log C$.
\begin{itemize}
\item \textbf{bound}: $\log C = \kappa + \log(2m(n-1))$, with one fitted constant $\kappa$.
\item \textbf{n,m}: OLS on $1,\log n,\log m$.
\item \textbf{n,m,diam}: OLS on $1,\log n,\log m,\log \mathrm{diam}$.
\item \textbf{n,m,diam,dmax}: OLS on $1,\log n,\log m,\log\mathrm{diam},\log\Delta$.
\end{itemize}
We evaluate each model by leave-one-family-out (LOFO) cross-validation: the model is fit on 14 families and predicts the held-out family. This asks whether the prediction transfers to a new graph type, which is the relevant question for a predictor that should be reliable beyond the training families. We report mean absolute log error (MAE), the median and 90th percentile of $|\hat C_{\rm pred}/\hat C-1|$, and a bootstrap 95\% CI of the MAE over graphs.

\section{Experiments}
Code and seeds are in the accompanying repository (see Reproducibility). The full run takes about 10 seconds on one CPU core.

\begin{table}[t]
\caption{Leave-one-family-out prediction error on 140 graphs.}
\label{tab:lofo}
\begin{center}
\begin{small}
\begin{tabular}{lcccc}
\toprule
Model & MAE (log) & Median $|$ratio$-1|$ & 90th pct & 95\% CI (MAE) \\
\midrule
bound & 0.937 & 0.649 & 3.15 & [0.83, 1.05] \\
n,m & 0.349 & 0.278 & 0.78 & [0.30, 0.40] \\
n,m,diam & 0.259 & 0.200 & 0.54 & [0.23, 0.29] \\
n,m,diam,dmax & \textbf{0.185} & \textbf{0.121} & \textbf{0.46} & [0.16, 0.22] \\
\bottomrule
\end{tabular}
\end{small}
\end{center}
\vskip -0.1in
\end{table}

Table~\ref{tab:lofo} shows that the four-statistic log-linear model reduces the LOFO error of the universal bound by a factor of about five, and that every added statistic helps. The fitted coefficients on all graphs are $\log C \approx -0.55 + 1.24\log n - 0.14\log m + 0.98\log\mathrm{diam} + 0.48\log\Delta$ (in-sample RMSE $0.19$ in log units). The large positive coefficient on $\log n$ together with a negative one on $\log m$ reflects that, for sparse graphs, $C$ scales with $n$ and diameter rather than with the edge count that dominates the bound.

Figure~\ref{fig:scatter} (left) shows that the bound is typically loose by a factor of about four and not uniformly so: complete graphs sit far below the line ($C/2m(n-1)\approx0.01$), while paths and random trees sit much closer. Figure~\ref{fig:per}, and the per-family table in the repository, show where the four-statistic model fails.

\begin{figure}[t]
\begin{center}
\includegraphics[width=\columnwidth]{../figs/per_family.pdf}
\end{center}
\caption{Median $|$prediction$/C-1|$ per family under LOFO for the four-statistic model. Families with 7 graphs or fewer are noisy.}
\label{fig:per}
\end{figure}

\paragraph{Where it works and where it fails.} Under LOFO, the four-statistic model has median error $\le 0.13$ on barbells, cycles, $G(n,0.1)$, $G(n,0.2)$, lollipops, paths, random trees, random 3-regular graphs, stars and wheels. It is poor on complete graphs (median error $0.52$), grids ($0.35$), and binary trees ($0.37$), and somewhat worse on dense $G(n,0.4)$ ($0.19$) and random 4-regular graphs ($0.18$). The failures share a pattern: the walk either mixes almost immediately (complete graphs, where $C\approx n\log n$ despite $m\approx n^2/2$) or has a geometry (grids, binary trees) whose cover time depends on more than the diameter and degrees.

\section{Discussion and Limitations}
The result is a practical, cheap predictor for small connected graphs of ``ordinary'' sparse types. Its accuracy depends on the family: for paths-like, tree-like and sparse random graphs an error of roughly 8--13\% is typical, but for dense or highly symmetric graphs it can be off by a factor of 1.5 or more. Using the bound as a baseline is misleading as soon as the graph is not a lollipop.

Limitations. (1) The test set is small and synthetic: 15 families with 7--21 graphs each, all with $n\le30$. Conclusions about families (e.g.\ ``grids fail'') rest on 7 graphs and should be treated as hypotheses. (2) Cover times are from uniformly random starts, not the worst start; worst-start values would be larger and would be better matched by the bound. (3) The predictors are fit with log-linear regression on four statistics; richer features (spectral gap, effective resistance, hitting-time sums) would likely do better but were not tested. (4) Our LOFO evaluation uses the family as the unit of generalisation. Extrapolation to graph types outside our generator, or to $n$ beyond 30, is not supported by the evidence here. (5) Monte Carlo error is small ($<1.2\%$), but we did not repeat the whole pipeline across independent graph-generation seeds, so variability from the choice of random graphs is not quantified beyond the three- or two-seed replicates within families.

\section{Reproducibility}
Code: \texttt{code/cover.py} (graph generation and simulation; default seed 0, 10{,}000 walks per graph via \texttt{--walks 10000}) and \texttt{code/analyze.py} (fits, LOFO evaluation, figures; bootstrap seed 0). Run from \texttt{code/}: \texttt{python cover.py --walks 10000 --seed 0 --out ../results/cover.csv} then \texttt{python analyze.py}. Dependencies: Python 3, numpy, networkx 3.x, matplotlib. Outputs: \texttt{results/cover.csv}, \texttt{results/summary.txt}, \texttt{results/lofo\_by\_family.csv}, \texttt{figs/*.pdf}.

\begin{figure}[t]
\begin{center}
\includegraphics[width=\columnwidth]{../figs/scatter.pdf}
\end{center}
\caption{Left: estimated cover time against $2m(n-1)$ for all 140 graphs, coloured by family; dashed line is equality. Right: LOFO predictions of the four-statistic model and the rescaled bound against $\hat C$.}
\label{fig:scatter}
\end{figure}

\bibliography{refs}
\bibliographystyle{icml2026}

\end{document}
